% !TeX root = ../../Book1.tex
% 纲要标题：### 16.3 Lie 子代数的积分与闭子群定理
% 纲要位置：几何学计划纲要.md，第 751 行。
% 正文状态：初稿；尚未正式审核。

% 纲要内容开始
% 
% * Lie 子代数的积分分布
% * 浸入子群的构造
% * 闭子群定理
% * 闭性与嵌入性的关系
% 
% 纲要内容结束

\section{Lie 子代数的积分与闭子群定理}
\label{g1:sec:1603}

单位元处的一组允许方向，怎样长成一个子群？线性子空间还不够，因为沿两种方向交替运动会产生括号方向；要求这些新方向仍被允许，就得到 Lie 子代数。Frobenius 定理将把它积分成经过单位元的最大叶，但这片叶首先只有自己的流形拓扑，未必采用原群的子空间拓扑。闭子群定理将进一步说明，若一个代数子群在原群中闭，单凭这个拓扑条件就足以恢复嵌入子流形结构。

\ProseParagraphBreak

以下 Lie 群均为有限维实、Hausdorff、第二可数且无边界。固定 Lie 群 \(G\)，单位元为 \(e\)，Lie 代数为 \(\mathfrak g=\mathrm{T}_eG\)。我们使用\cref{g1:thm:frobenius,g1:thm:maximal-connected-leaf}中已经构造的最大积分叶，并使用上一节的指数映射及局部逆。本节只讨论子群自身的流形结构；陪集空间的流形结构将在下一章构造。

\subsection{Lie 子代数的积分分布}
\label{g1:subsec:1603:01}

回顾前面的定义：Lie 子代数是满足 \([\mathfrak h,\mathfrak h]\subseteq\mathfrak h\) 的线性子空间 \(\mathfrak h\subseteq\mathfrak g\)。这个条件可以在整个群上展开：把 \(\mathfrak h\) 从单位元左移到每个点，得到
\begin{equation}\label{g1:eq:subalgebra-left-distribution}
	D_g=(\mathrm{d}L_g)_e\mathfrak h,\qquad g\in G.
\end{equation}
这里 \(L_g(a)=ga\)，所以每个 \(D_g\) 都有相同维数。我们将证明括号封闭性正好保证这个分布可积。

\begin{proposition}[左不变积分分布]\label{g1:prop:lie-subalgebra-left-distribution}
	设 \(G\) 为有限维实 Lie 群，\(\mathfrak g=\mathrm{T}_eG\)，\(\mathfrak h\subseteq\mathfrak g\) 为 \(r\) 维 Lie 子代数。令 \(D_g=(\mathrm{d}L_g)_e\mathfrak h\)，即\eqref{g1:eq:subalgebra-left-distribution}中的分布。则 \(D\) 是秩 \(r\) 的光滑对合分布，所有左平移都保持此分布。对 \(X\in\mathfrak h\)，还有
	\[
		\operatorname{Ad}_{\exp(tX)}\mathfrak h=\mathfrak h\qquad(t\in\mathbb R).
	\]
\end{proposition}
\begin{proof}
	选 \(\mathfrak h\) 的基 \(X_1,\ldots,X_r\)，将其延成左不变向量场 \(X_1^L,\ldots,X_r^L\)。这些场光滑，且在每点线性无关并张成 \(D_g\)，故给出光滑分布。Lie 代数括号的定义及 \(\mathfrak h\) 的封闭性说明，每个 \([X_i^L,X_j^L]\) 都是这些基场的常系数线性组合。对一般局部截面 \(U=\sum_i f_iX_i^L\)、\(V=\sum_j h_jX_j^L\)，有
	\[
		[U,V]=\sum_{i,j}\left(f_ih_j[X_i^L,X_j^L]
		+f_iX_i^L(h_j)X_j^L-h_jX_j^L(f_i)X_i^L\right),
	\]
	仍取值于 \(D\)，所以分布对合。又因为 \(L_aL_g=L_{ag}\)，有 \(\mathrm{d}L_a(D_g)=D_{ag}\)，证明左平移不变性。

	最后固定 \(X,Y\in\mathfrak h\)。由\cref{g1:prop:exponential-naturality}，伴随作用与算子指数满足
	\[
		\operatorname{Ad}_{\exp(tX)}Y
		=e^{t\operatorname{ad}_X}Y
		=\sum_{j=0}^{\infty}\frac{t^j}{j!}\operatorname{ad}_X^jY.
	\]
	括号封闭性说明 \(\operatorname{ad}_X=[X,\cdot]\) 保持 \(\mathfrak h\)，所以级数的每个部分和均在 \(\mathfrak h\) 中。有限维线性子空间闭，级数极限也在其中。因此 \(\operatorname{Ad}_{\exp(tX)}\mathfrak h\subseteq\mathfrak h\)；把 \(t\) 换成 \(-t\)，再用两伴随算子互逆，得到反向包含。
\end{proof}

Frobenius 定理现在给出叶状坐标盒，在其中 \(D\) 由前 \(r\) 个坐标方向张成。后面的横向坐标在每张叶片上恒定。我们会反复使用这个局部事实来判断一个环境光滑映射何时也对叶自身的拓扑光滑。

\begin{lemma}[取值于叶的光滑映射]\label{g1:lem:leaf-valued-map}
	设 \(M,N\) 为有限维、Hausdorff、第二可数的无边界光滑流形，\(D\) 为 \(M\) 上的常秩光滑对合分布。若 \(N\) 连通，\(F:N\to M\) 光滑，且
	\[
		\mathrm{d}F_q(\mathrm{T}_qN)\subseteq D_{F(q)}
	\]
	对每个 \(q\in N\) 成立，则 \(F(N)\) 包含在一片最大叶中。将值域限制为该叶后，\(F\) 对叶的光滑结构仍然光滑。
\end{lemma}
\begin{proof}
	固定 \(q\in N\)，取连通的小坐标邻域 \(V\ni q\)，使 \(F(V)\) 落在一个叶状坐标盒中，盒内坐标记为 \((u,v)\)，其中 \(v\) 是横向坐标。由于 \(\mathrm{d}v\) 在 \(D\) 上为零，链式法则给出 \(\mathrm{d}(v\circ F)=0\)。在连通坐标小球内沿线段积分导数，便知 \(v\circ F\) 恒为一个常值。因此 \(F(V)\) 位于一张叶片中。

	在 \(N\) 中，把像位于同一最大叶的点归成一类。上面的局部结论说明这些类都是开集，它们两两不交并覆盖 \(N\)，故连通性使只有一个非空类。这证明全像位于一片叶。每个 \(F(V)\) 又位于该叶的一张开叶片中，叶片坐标下的 \(F\) 正是光滑函数 \(u\circ F\)，所以到叶的映射局部光滑，进而整体光滑。这里没有要求 \(F\) 是浸入，也没有把叶拓扑替换成环境子空间拓扑。
\end{proof}

\subsection{浸入子群的构造}
\label{g1:subsec:1603:02}

先说清要构造哪一种子群。一个子群 \(H\subseteq G\) 可以带有自己的 Lie 群结构；只要求包含映射光滑且微分单射时，得到的是浸入子群。如果还要求这份拓扑等于 \(G\) 在 \(H\) 上诱导的子空间拓扑，才得到嵌入子群。两种情形中的群乘法都必须对所用的流形结构光滑，不能从集合上的子群性质直接略过这一步。

\begin{definition}[浸入 Lie 子群]\label{g1:def:immersed-lie-subgroup}
	设 \(G\) 为有限维实 Lie 群。子群 \(H\subseteq G\) 若带有有限维、Hausdorff、第二可数的无边界 Lie 群结构，且包含映射 \(i:H\hookrightarrow G\) 为光滑浸入，则称为\term[jinruliziqun]{浸入 Lie 子群}[immersed Lie subgroup]。若包含映射还是光滑嵌入，则称为嵌入 Lie 子群。通过 \(\mathrm{d}i_e\) 将 \(\mathrm{T}_eH\) 识别为 \(\mathfrak g\) 的子空间。
\end{definition}

\begin{theorem}[Lie 子代数的积分]\label{g1:thm:lie-subalgebra-integration}
	设 \(G\) 为有限维实、Hausdorff、第二可数的无边界 Lie 群，\(\mathfrak h\subseteq\mathfrak g\) 为 Lie 子代数。存在唯一的连通浸入 Lie 子群 \(H\subseteq G\)，满足 \(\operatorname{Lie}(H)=\mathfrak h\)。其底集合及流形结构就是左不变分布 \(D_g=(\mathrm{d}L_g)_e\mathfrak h\) 经过 \(e\) 的最大叶，并且
	\begin{equation}\label{g1:eq:integral-subgroup-generated-exponentials}
		H=\langle\exp_G(\mathfrak h)\rangle.
	\end{equation}
	这里右侧表示由全部 \(\exp_G X\)、\(X\in\mathfrak h\) 生成的代数子群。
\end{theorem}
\begin{proof}
	由\cref{g1:prop:lie-subalgebra-left-distribution,g1:thm:maximal-connected-leaf}，经过 \(e\) 的最大叶 \(H\) 是连通 \(r\) 维流形，包含映射为浸入，且 \(\mathrm{T}_gH\) 在 \(\mathrm{T}_gG\) 中的像是 \(D_g\)。先提醒叶拓扑所需的两个可数与分离条件。它使每张局部叶片为开坐标片，包含映射连续；环境中的两个不交开邻域限制到叶上仍然开，所以 \(H\) Hausdorff。第二可数性则来自叶片构造：在 \(G\) 的可数个叶状坐标盒中，一张叶片与任意另一盒的交只有可数个连通分支，每个分支只落在那一盒的一张叶片中。沿有限相交链逐步加入叶片，每一步至多可数，故经过 \(e\) 的叶只需可数张叶片覆盖。取各片可数坐标基之并，得到叶的可数基。这正是最大叶定理中用于叶拓扑的论证，并非从环境子空间拓扑直接继承。

	左平移 \(L_a\) 保持 \(D\)，所以它把经过 \(e\) 的最大叶送到经过 \(a\) 的最大叶：若后者还有更大的连通积分流形，左移回来就会违背前者的最大性。若 \(a\in H\)，经过 \(a\) 的叶仍是 \(H\)，故 \(aH=H\)。于是 \(a,b\in H\) 蕴含 \(ab\in H\)，而 \(e\in aH\) 蕴含 \(a^{-1}\in H\)。这先证明 \(H\) 在集合上是子群。各 \(L_a|_H\)、\(a\in H\)，由\cref{g1:lem:leaf-valued-map}光滑，逆 \(L_{a^{-1}}|_H\) 同样光滑，所以已经是叶流形上的微分同胚。

	若 \(X\in\mathfrak h\)，曲线 \(t\mapsto\exp(tX)\) 的速度为 \(X^L\)，始终切于 \(D\)。它经过 \(e\)，故由取值于叶的映射引理，整条曲线在 \(H\) 中且对叶拓扑光滑。令 \(H'\) 是由这些曲线的值生成的子群，则 \(H'\subseteq H\)。为证明相等，选 \(\mathfrak h\) 的基 \(X_1,\ldots,X_r\)，考虑
	\[
		f(t^1,\ldots,t^r)=\exp(t^1X_1)\cdots\exp(t^rX_r).
	\]
	其值在 \(H'\) 中。记第 \(i\) 个因子之后的乘积为 \(b_i\)。对 \(t^i\) 求导，再用左平移送回单位元，所得方向为 \(\operatorname{Ad}_{b_i^{-1}}X_i\)。每个 \(b_i\) 都是 \(\exp(\mathfrak h)\) 中元素的乘积，上一命题说明它的伴随作用保持 \(\mathfrak h\)，故 \(\mathrm{d}f\) 的像包含于 \(D\)。取值于叶的映射引理使 \(f:\mathbb R^r\to H\) 光滑。在零点，它的微分把标准基送到 \(X_1,\ldots,X_r\)，所以到 \(\mathrm{T}_eH\) 为同构。逆函数定理说明 \(H'\) 包含 \(H\) 中单位元的一个开邻域。

	由 \(H\) 中的左平移已是微分同胚，\(H'\) 是 \(H\) 中的开子群；它的其他左陪集也开，故补集开。连通性迫使 \(H'=H\)，证明\eqref{g1:eq:integral-subgroup-generated-exponentials}。同时，生成元的伴随作用保持 \(\mathfrak h\)，于是
	\[
		\operatorname{Ad}_a\mathfrak h=\mathfrak h\qquad(a\in H).
	\]

	现在验证群运算在叶拓扑下光滑。对 \(a,b\in H\)、\(X,Y\in\mathfrak h\)，把切向量写成左平移的形式，乘法 \(\mu\) 的微分为
	\begin{equation}\label{g1:eq:subgroup-multiplication-differential}
		\mathrm{d}\mu_{(a,b)}\bigl(\mathrm{d}L_aX,\mathrm{d}L_bY\bigr)
		=\mathrm{d}L_{ab}\bigl(\operatorname{Ad}_{b^{-1}}X+Y\bigr)\in D_{ab}.
	\end{equation}
	这个公式可对曲线
	\[
		a\exp(tX)b\exp(tY)
		=ab\,\exp\bigl(t\operatorname{Ad}_{b^{-1}}X\bigr)\exp(tY)
	\]
	在零点求导得到。环境乘法组成的映射 \(H\times H\to G\) 光滑，微分又取值于 \(D\)，所以取值于叶的映射引理表明它作为 \(H\times H\to H\) 也光滑。对逆映射 \(\iota(a)=a^{-1}\)，同理有
	\[
		\mathrm{d}\iota_a(\mathrm{d}L_aX)
		=\mathrm{d}L_{a^{-1}}(-\operatorname{Ad}_aX)\in D_{a^{-1}},
	\]
	这来自 \((a\exp(tX))^{-1}=a^{-1}\exp(-t\operatorname{Ad}_aX)\)。因此逆映射也对叶结构光滑，\(H\) 确实成为 Lie 群。

	包含映射是光滑群同态，单位元处微分像为 \(\mathfrak h\)，所以其 Lie 代数恰为给定子代数。若 \(K\subseteq G\) 是另一个连通浸入 Lie 子群，且单位元切空间同样为 \(\mathfrak h\)，同态与左平移相容使它在每点的切空间像为 \(D\)。因此 \(K\) 是经过 \(e\) 的连通积分流形。最大叶定理给出包含 \(K\to H\)，并保证它为开嵌入。其像又是 \(H\) 的子群，故是开子群；由 \(H\) 连通，像等于 \(H\)。这个开嵌入遂为微分同胚，证明底集合和 Lie 群结构都唯一。若 \(r=0\)，最大叶只是 \(\{e\}\)，上述结论同样成立。
\end{proof}

浸入子群使用叶拓扑，群运算的光滑性由上面的三个性质保证：微分切于分布、横向坐标恒定，以及叶片本身是叶中的开集。对嵌入子群，则可直接调用\cref{g1:prop:embedded-initial-property}中的限制映射性质。

\subsection{闭子群定理}
\label{g1:subsec:1603:03}

现在反向开始：只给 \(G\) 的一个代数子群 \(K\)，暂时不赋予任何流形结构。若它闭，我们可以用全部留在 \(K\) 内的单参数子群识别候选切空间，再证明附近没有额外的横向部分。闭性将在两次取极限时起决定作用。

\begin{lemma}[小元素的极限方向]\label{g1:lem:closed-subgroup-limit-direction}
	设 \(K\) 是有限维实 Lie 群 \(G\) 的闭子群，给 \(\mathfrak g\) 任选 Euclid 范数。若 \(Z_j\ne0\)、\(Z_j\to0\)、\(\exp Z_j\in K\)，且 \(Z_j/\lVert Z_j\rVert\to Z\)，则
	\[
		\exp(tZ)\in K\qquad\text{对每个 }t\in\mathbb R.
	\]
\end{lemma}
\begin{proof}
	固定 \(t\in\mathbb R\)，令 \(\varepsilon_j=\lVert Z_j\rVert\)，选整数 \(m_j\) 使 \(|m_j\varepsilon_j-t|\le\varepsilon_j\)。因为 \(\varepsilon_j\to0\)，有
	\[
		m_jZ_j=(m_j\varepsilon_j)\frac{Z_j}{\varepsilon_j}\longrightarrow tZ.
	\]
	单参数子群规律给出
	\[
		(\exp Z_j)^{m_j}=\exp(m_jZ_j)\in K.
	\]
	指数映射连续，而 \(K\) 闭，故极限 \(\exp(tZ)\) 也属于 \(K\)。整数 \(m_j\) 可以为负，所以论证同时包含正、负时间；\(t\) 任意，结论成立。
\end{proof}

\begin{theorem}[闭子群定理]\label{g1:thm:closed-subgroup}
	设 \(G\) 为有限维实、Hausdorff、第二可数的无边界 Lie 群，\(K\subseteq G\) 为闭的代数子群。则 \(K\) 具有唯一的嵌入 Lie 子群结构，其 Lie 代数为
	\begin{equation}\label{g1:eq:closed-subgroup-lie-algebra}
		\mathfrak k=\{X\in\mathfrak g:\exp(tX)\in K\text{ 对所有 }t\in\mathbb R\}.
	\end{equation}
	对任意线性补空间 \(\mathfrak g=\mathfrak k\oplus\mathfrak m\)，可以选择零点的邻域 \(V\subseteq\mathfrak k\)、\(W\subseteq\mathfrak m\)，使
	\[
		\Psi:V\times W\longrightarrow U,\qquad \Psi(X,Y)=\exp X\exp Y
	\]
	为到 \(G\) 中单位元开邻域的微分同胚，并且
	\begin{equation}\label{g1:eq:closed-subgroup-local-slice}
		K\cap U=\Psi(V\times\{0\})=\exp V.
	\end{equation}
\end{theorem}
\begin{proof}
	先证明\eqref{g1:eq:closed-subgroup-lie-algebra}确实给出 Lie 子代数。对每个固定 \(t\)，\(X\mapsto\exp(tX)\) 连续，所以 \(\mathfrak k\) 是闭集 \(K\) 的这些原像之交，因而在 \(\mathfrak g\) 中闭。它包含零，并且对任意实数倍封闭，因为 \(t(cX)=(tc)X\)。若 \(X,Y\in\mathfrak k\)，由\cref{g1:cor:lie-product-limit}，对每个实数 \(t\) 有
	\[
		\exp(t(X+Y))
		=\lim_{n\to\infty}\bigl(\exp(tX/n)\exp(tY/n)\bigr)^n.
	\]
	右侧每一项都在子群 \(K\) 中，闭性使极限也在 \(K\) 中。因此 \(X+Y\in\mathfrak k\)，\(\mathfrak k\) 是闭线性子空间。

	对 \(k\in K\)，\cref{g1:prop:exponential-naturality}说明指数与共轭相容，故
	\[
		\exp(t\operatorname{Ad}_kX)=k\exp(tX)k^{-1}\in K
		\qquad(X\in\mathfrak k).
	\]
	于是 \(\operatorname{Ad}_k\mathfrak k=\mathfrak k\)，反向包含使用 \(k^{-1}\)。若 \(X,Y\in\mathfrak k\)，则 \(\exp(sX)\in K\)，所以 \(\operatorname{Ad}_{\exp(sX)}Y\) 始终位于闭线性空间 \(\mathfrak k\)。在零点取差商极限，得到
	\[
		[X,Y]=\left.\frac{\mathrm{d}}{\mathrm{d}s}\right|_{s=0}
		\operatorname{Ad}_{\exp(sX)}Y\in\mathfrak k.
	\]
	这证明其括号封闭性。

	取补空间 \(\mathfrak m\)。我们先证明存在 \(\delta>0\)，使
	\[
		Y\in\mathfrak m,\quad \lVert Y\rVert<\delta,\quad \exp Y\in K
		\quad\Longrightarrow\quad Y=0.
	\]
	若不存在，就可取非零 \(Y_j\in\mathfrak m\)，满足 \(Y_j\to0\)、\(\exp Y_j\in K\)。\(\mathfrak m\) 的单位球面紧，故取子列后 \(Y_j/\lVert Y_j\rVert\to Y\in\mathfrak m\)，且 \(\lVert Y\rVert=1\)。由\cref{g1:lem:closed-subgroup-limit-direction}，\(\exp(tY)\in K\) 对所有 \(t\) 成立，即 \(Y\in\mathfrak k\)。这与 \(\mathfrak k\cap\mathfrak m=\{0\}\) 矛盾。若 \(\mathfrak m=\{0\}\)，这一断言直接成立。

	考虑光滑映射 \(\Psi(X,Y)=\exp X\exp Y\)。\cref{g1:prop:exponential-local-diffeomorphism}给出 \(\mathrm{d}(\exp)_0=\operatorname{id}_{\mathfrak g}\)，群乘法在 \((e,e)\) 的微分是向量相加，故
	\[
		\mathrm{d}\Psi_{(0,0)}(X,Y)=X+Y.
	\]
	直和分解使它可逆。逆函数定理给出零点附近的局部微分同胚；再把定义域缩成乘积 \(V\times W\)，并使 \(W\) 包含于上面的 \(\delta\) 球中。这个乘积的像 \(U\) 仍在 \(G\) 中开。

	若 \(k\in K\cap U\)，唯一写成 \(k=\exp X\exp Y\)，其中 \(X\in V\)、\(Y\in W\)。由于 \(X\in\mathfrak k\)，有 \(\exp X\in K\)，所以
	\[
		\exp Y=(\exp X)^{-1}k\in K.
	\]
	横向排除结论强制 \(Y=0\)。反过来，\(X\in V\) 时 \(\exp X\in K\cap U\)。这就证明了完整的切片等式\eqref{g1:eq:closed-subgroup-local-slice}；它排除了附近所有额外分支，而不仅仅给出一个包含于 \(K\) 的参数片。

	将这张切片用 \(L_k\)、\(k\in K\)，移到每个子群点附近，得到 \(K\) 的嵌入子流形结构。它取环境子空间拓扑，Hausdorff 性和第二可数性从 \(G\) 继承；切片坐标转换由环境坐标转换限制得到，故光滑。群乘法与逆映射在 \(G\) 上光滑，其限制值域在 \(K\) 中，\cref{g1:prop:embedded-initial-property}因而保证它们在这个子流形结构下光滑。所以 \(K\) 是嵌入 Lie 子群。单位元的切片为 \(X\mapsto\exp X\)、\(X\in V\)，其微分为包含 \(\mathfrak k\hookrightarrow\mathfrak g\)，故 Lie 代数恰为 \(\mathfrak k\)。最后，嵌入子流形的光滑结构由包含映射的初始性质唯一决定，证明唯一性。
\end{proof}

当 \(\mathfrak k=\{0\}\) 时，局部切片只是单位元，左移以后每个点都孤立，所以 \(K\) 是离散子群。当 \(\mathfrak k=\mathfrak g\) 时，切片占满单位元的一个开邻域，所以 \(K\) 是开子群。两种极端都包含在同一证明中。

\ProseParagraphBreak

Lie 代数只描述单位元所在的连通部分，不能确定全部离散分支。例如加法群 \(\mathbb R\) 中，\(\{0\}\) 与 \(\mathbb Z\) 都是闭子群，二者 Lie 代数都为零，但它们不是同一个子群。因此前面的积分唯一性明确要求连通。对一般闭子群 \(K\)，单位元连通分支 \(K^0\) 是开连通子群：流形局部路径连通使连通分支开，乘法与逆映射又把含单位元的连通集送入同一分支。它具有与 \(K\) 相同的单位元切空间，所以正是 \(\mathfrak k\) 的唯一连通积分子群。

\subsection{闭性与嵌入性的关系}
\label{g1:subsec:1603:04}

一般嵌入子流形不必闭，例如平面中的开线段。然而子群的平移结构会把单位元附近的局部闭性传播到整个群。另一方面，如果一个已有自身拓扑的浸入 Lie 子群碰巧具有闭像，还须说明原拓扑与闭子群定理给出的拓扑一致，才能说原包含映射是嵌入。

\begin{proposition}[闭性与嵌入性]\label{g1:prop:immersed-subgroup-closed-embedded}
	设 \(H\) 为有限维实 Lie 群 \(G\) 的浸入 Lie 子群，两个 Lie 群都为 Hausdorff、第二可数。则包含映射 \(H\hookrightarrow G\) 是嵌入，当且仅当子集 \(H\) 在 \(G\) 中闭。
\end{proposition}
\begin{proof}
	先设包含映射是嵌入。由\cref{g1:prop:embedding-slice-characterization}，存在 \(G\) 中单位元的开邻域 \(U\)，使 \(H\cap U\) 是一张坐标切片，因而在 \(U\) 中闭。令 \(K=\overline H\)。乘法与逆映射连续，故 \(K\) 仍是子群：逼近 \(a,b\in K\) 的子群点列，其乘积及逆分别逼近 \(ab,a^{-1}\)，这些极限仍在闭包中。因为 \(U\) 开，局部闭性给出 \(K\cap U=H\cap U\)。于是 \(H\) 包含 \(K\) 中单位元的开邻域，故是 \(K\) 的开子群。每个左陪集也开，所以 \(H\) 的补集在 \(K\) 中开，即 \(H\) 又在 \(K\) 中闭。它在 \(K=\overline H\) 中稠密，只能有 \(H=K\)，证明闭性。

	反过来，假设像 \(H\) 在 \(G\) 中闭，令 \(K\) 表示同一个集合上由闭子群定理给出的嵌入 Lie 群结构。原包含映射通过 \(K\) 分解为一个双射光滑群同态 \(j:H\to K\)，光滑性使用嵌入子流形的初始性质。需要证明 \(j^{-1}\) 连续，而不是把这一点当作双射的自动结论。

	给定 \(H\) 中单位元的任意开邻域 \(O\)，可取较小开邻域 \(V\)，使 \(C=\overline V\) 紧且 \(C^{-1}C\subseteq O\)。具体可先由乘法与逆的连续性取 \(W\) 满足 \(W^{-1}W\subseteq O\)，再在坐标图中取闭包包含于 \(W\) 的小球作为 \(V\)。第二可数性使覆盖 \(\{hV:h\in H\}\) 有可数子覆盖，故对某列 \(h_\ell\) 有
	\[
		K=\bigcup_{\ell=1}^{\infty}j(h_\ell)\,j(C).
	\]
	每个 \(j(C)\) 都紧，在 Hausdorff 流形 \(K\) 中闭。它不能没有内点。否则右侧是可数个闭的无内点集；在 \(K\) 的一张坐标图中先取一个具有非空内部的紧闭球，再逐次在前一闭球的内部选一个更小的非空闭球，避开第 \(\ell\) 个闭集。这些嵌套紧集有公共点，而该点不属于任何一个所列闭集，与覆盖矛盾。零维时闭的无内点集只能为空，结论直接成立。

	因此可取 \(K\) 中非空开集 \(W_0\subseteq j(C)\) 及 \(k_0\in W_0\)。则
	\[
		k_0^{-1}W_0\subseteq j(C)^{-1}j(C)
		=j(C^{-1}C)\subseteq j(O).
	\]
	左侧是单位元的开邻域，所以 \(j(O)\) 包含单位元邻域。在 \(O\) 的每个点先平移一个更小的单位元邻域，并应用相同论证，可知 \(j(O)\) 本身开。再用群平移，\(j\) 把任意开集送到开集。于是双射 \(j\) 为同胚，原 \(H\) 的拓扑等于 \(G\) 的子空间拓扑。原包含映射又已知是浸入，故按定义为光滑嵌入。
\end{proof}

这个证明中的第二可数性用于取得可数紧集覆盖；紧性和环境的 Hausdorff 性用于把它们的像看成闭集。它们共同排除了闭像还带有一套更细 Lie 群拓扑的可能。对于一般的单射浸入映射，没有群平移结构时，不能照搬这个结论。

\begin{example}[无理斜率子群]\label{g1:exm:irrational-torus-subgroup}
	在环面 Lie 群 \(T^2=\mathbb R^2/\mathbb Z^2\) 中，固定无理数 \(\alpha\)，定义
	\[
		f_\alpha:\mathbb R\longrightarrow T^2,\qquad
		f_\alpha(t)=[t,\alpha t].
	\]
	它是光滑群同态，微分像为直线 \(\mathfrak h_\alpha=\mathbb R(1,\alpha)\)。若 \(f_\alpha(t)=0\)，则 \(t\in\mathbb Z\) 且 \(\alpha t\in\mathbb Z\)，无理性迫使 \(t=0\)，所以它单射；微分处处非零，故为浸入。将像 \(H_\alpha\) 赋予从 \(\mathbb R\) 传来的 Lie 群结构，便得到 \(\mathfrak h_\alpha\) 的连通积分子群。

	\cref{g1:prop:linear-torus-leaves}已经证明这条轨道在环面中稠密。其关键是整数倍 \(\alpha\) 模 \(1\) 稠密：对任意指定的第一坐标 \(u\)，时刻 \(u+n\) 的第二坐标为 \(\alpha u+n\alpha\)，因此能够进入任意给定小矩形。这个像却不是整个环面，因为它与第一坐标为零的圆周只相交于可数个点 \([0,n\alpha]\)。于是 \(H_\alpha\) 是真稠密子群，既不闭也不嵌入。

	拓扑差异还可直接检测。前述整数逼近给出 \(|n_j|\to\infty\)，但 \(n_j\alpha\bmod1\to0\)，所以 \(f_\alpha(n_j)\to f_\alpha(0)\) 在环面中成立；在 \(H_\alpha\) 自身的参数坐标里，\(n_j\) 却不趋于零。这正是从子空间像返回叶参数的逆映射不连续。子群自身始终是一维 Lie 群，闭包却是二维环面，其 Lie 代数也从一条直线增大为整个 \(\mathbb R^2\)。
\end{example}

因此，Lie 子代数的积分定理提供连通浸入子群，闭子群定理提供采用子空间拓扑的嵌入子群，两者处理的条件不同。后面构造齐性空间时，将用闭性把子群附近的横向坐标组织为陪集空间的光滑坐标。

% 纲要条目排为编号小节标题，并进入目录与书签；说明段仍保留于纲要注释中。
